\documentclass[10pt,a4paper]{article}
\usepackage[utf8]{inputenc}
\usepackage[T1]{fontenc}
\usepackage[portuguese]{babel}
\usepackage[margin=1.2cm,top=1.2cm,bottom=1.2cm]{geometry}
\usepackage{amsmath,amssymb,amsfonts}
\usepackage{enumitem}

\pagestyle{empty}

\begin{document}

% --- CABEÇALHO COMPACTO ---
\noindent
\textbf{UFSC --- FSC5101 Física I --- Prova 1 (Capítulos 1, 2 e 3)} \hfill Nome: \underline{\hspace{5.0cm}} Matrícula: \underline{\hspace{2.2cm}} Data: \underline{\hspace{1.3cm}} Nota: \underline{\hspace{1.0cm}}

\vspace{2mm}
\noindent
{\small \textit{\textbf{Instruções:} Responda a prova a caneta. Todas as 5 questões valem 2,0 pontos cada. Adote $g = 9,80\text{ m/s}^2$. É obrigatório utilizar a notação vetorial apropriada (ex: $\vec{v}$ ou $\hat{i}, \hat{j}, \hat{k}$) e indicar explicitamente todas as unidades nas respostas.}}

\vspace{2mm}
\hrule
\vspace{4mm}

% --- QUESTÃO 1 ---
\noindent\textbf{1.} Dados os vetores no espaço tridimensional $xyz$:
$$\vec{A} = 6,0\,\hat{i} - 2,0\,\hat{j} + 3,0\,\hat{k}\text{ (m)} \quad \text{e} \quad \vec{B} = 4,0\,\hat{i} + 4,0\,\hat{j} - 2,0\,\hat{k}\text{ (m)}$$
Determine:
\begin{enumerate}[label=\alph*),itemsep=2pt,topsep=3pt]
    \item Os módulos $|\vec{A}|$, $|\vec{B}|$ e o produto escalar $\vec{A} \cdot \vec{B}$.
    \item O ângulo $\theta$ formado entre os vetores $\vec{A}$ e $\vec{B}$.
    \item O produto vetorial $\vec{C} = \vec{A} \times \vec{B}$ em termos dos vetores unitários $(\hat{i}, \hat{j}, \hat{k})$.
\end{enumerate}

\vfill
\hrule
\vspace{4mm}

% --- QUESTÃO 2 ---
\noindent\textbf{2.} Um trem de passageiros trafega em uma via retilínea a uma velocidade constante de $v_0 = 108\text{ km/h}$ ($30,0\text{ m/s}$). Repentinamente, o maquinista avista um trecho em manutenção a uma distância de $D = 180,0\text{ m}$ à sua frente. O maquinista possui um tempo de reação de $t_{\text{reac}} = 0,60\text{ s}$ antes de acionar efetivamente o sistema de freios de emergência. Determine:
\begin{enumerate}[label=\alph*),itemsep=2pt,topsep=3pt]
    \item A distância percorrida pelo trem durante o tempo de reação e a distância restante disponível para a frenagem.
    \item O módulo da desaceleração constante mínima $a$ (em $\text{m/s}^2$) necessária para que o trem pare exatamente $12,0\text{ m}$ antes de atingir o trecho em manutenção.
\end{enumerate}

\vfill

\newpage

% --- QUESTÃO 3 ---
\noindent\textbf{3.} Do topo de um edifício de altura $H = 122,5\text{ m}$, uma pedra A é solta a partir do repouso no instante $t = 0$. Exatamente $\Delta t = 1,00\text{ s}$ após a soltura da pedra A, uma segunda pedra B é lançada verticalmente para baixo a partir do mesmo ponto inicial. Adotando $g = 9,80\text{ m/s}^2$, determine:
\begin{enumerate}[label=\alph*),itemsep=2pt,topsep=3pt]
    \item O tempo total $t_A$ necessário para a pedra A atingir o solo.
    \item O módulo da velocidade inicial $v_{0B}$ com que a pedra B deve ser lançada para que ambas as pedras atinjam o solo no mesmo instante exato.
\end{enumerate}

\vfill
\hrule
\vspace{4mm}

% --- QUESTÃO 4 ---
\noindent\textbf{4.} Um avião voa em relação ao ar com uma velocidade constante de módulo $v_{A/R} = 120,0\text{ km/h}$ direcionado para a direção Leste. Contudo, sobra um vento constante vindo do Sul em direção ao Norte com velocidade em relação à Terra de $v_{R/T} = 50,0\text{ km/h}$. Determine:
\begin{enumerate}[label=\alph*),itemsep=2pt,topsep=3pt]
    \item O módulo da velocidade do avião $v_{A/T}$ em relação a um observador parado na superfície da Terra.
    \item O ângulo de desvio $\theta$ da trajetória resultante do avião em relação à direção Leste.
\end{enumerate}

\vfill

\newpage

% --- QUESTÃO 5 ---
\noindent\textbf{5.} Um projétil é disparado a partir do solo plano e horizontal com uma velocidade inicial $v_0 = 50,0\text{ m/s}$ sob um ângulo de elevação de $\theta_0 = 36,9^\circ$ ($\sin 36,9^\circ = 0,60$; $\cos 36,9^\circ = 0,80$). O projétil atinge o topo de uma torre vertical exatamente no instante $t = 5,00\text{ s}$ após o disparo. Adotando $g = 9,80\text{ m/s}^2$, determine:
\begin{enumerate}[label=\alph*),itemsep=2pt,topsep=3pt]
    \item A altura $H$ da torre e a distância horizontal $D$ do ponto de disparo até a base da torre.
    \item O módulo $v$ e a direção (ângulo em relação à horizontal) do vetor velocidade do projétil no instante exato do impacto com o topo da torre.
\end{enumerate}

\vfill

\end{document}
